\begin{statementbox}
\textbf{C089\quad 双曲线切线与两条渐近线围成三角形的面积定值}

\textbf{对象与记号。}设
\[
H:\ \frac{x^2}{a^2}-\frac{y^2}{b^2}=1\quad (a,b>0),\qquad
\ell_1:y=\frac ba x,\quad \ell_2:y=-\frac ba x.
\]
$O=(0,0)$ 是两条渐近线的交点（双曲线中心）。任选切点
$P=(x_0,y_0)\in H$，在 $P$ 处的切线 $t_P$ 分别与
$\ell_1,\ell_2$ 交于 $A,B$。\textbf{$A$ 固定在正斜率渐近线上，
$B$ 固定在负斜率渐近线上}，以免后续公式混淆。

\textbf{主结论（恒等式）：}
\[
\boxed{S_{\triangle AOB}=ab}\qquad\text{与切点 $P$ 无关}.
\]
\begin{center}
\begin{tikzpicture}[x=1.12cm,y=1.12cm,>=stealth,every node/.style={font=\small}]
\coordinate (O) at (0,0);
\coordinate (A) at (3.24,1.98);
\coordinate (B) at (1,-0.611111);
\coordinate (P) at (2.12,0.684444);
\fill[red!7] (O)--(A)--(B)--cycle;
\draw[gray!65,->] (-0.65,0)--(4.15,0) node[right] {$x$};
\draw[gray!65,->] (0,-1.35)--(0,2.6) node[above] {$y$};
\draw[blue!65!black,densely dashed] (-0.65,-0.3972)--(3.8,2.3222);
\draw[blue!65!black,densely dashed] (-0.65,0.3972)--(2.05,-1.2528);
\draw[teal!70!black,thick,domain=-0.9:1.22,samples=90,variable=\t]
 plot ({1.8/2*(exp(\t)+exp(-\t))},{1.1/2*(exp(\t)-exp(-\t))});
\draw[red!75!black,very thick] ($(B)!-0.15!(A)$)--($(A)!-0.12!(B)$);
\draw[orange!90!black,thick] (A)--(O)--(B);
\fill (O) circle(1.5pt) node[below left=5pt] {$O$};
\fill (A) circle(1.5pt) node[above left] {$A$};
\fill (B) circle(1.5pt) node[below=9pt] {$B$};
\fill (P) circle(1.5pt) node[left=4pt] {$P$};

\node[anchor=west] at (4.55,2.05) {\textcolor{teal!70!black}{双曲线 $H$}};
\node[anchor=west] at (4.55,1.35) {\textcolor{red!75!black}{切线 $t_P$}};
\node[anchor=west] at (4.55,0.65) {$A\in\ell_1,\ B\in\ell_2$};
\node[anchor=west] at (4.55,-0.05) {$S_{\triangle AOB}=ab$};
\node[anchor=west] at (4.55,-0.75) {$P\text{ 是 }AB\text{ 中点}$};
\end{tikzpicture}

\end{center}

\textbf{伴随性质（同样对每个切点成立）：}
\[
\boxed{P=\frac{A+B}{2},\quad PA=PB},\qquad
\boxed{OA\cdot OB=a^2+b^2=c^2},\quad c=\sqrt{a^2+b^2}.
\]
因此，以 $OP$ 为中线，三角形被平分：
\[
\boxed{S_{\triangle OAP}=S_{\triangle OPB}=\frac{ab}{2}}.
\]
\textbf{成立范围：}右支、左支、上下方的任意有限切点均成立；
顶点 $P=(\pm a,0)$ 的竖直切线 $x=\pm a$ 也包括在内。
切线与两条渐近线分别有唯一交点 $A,B$，且 $A\neq B$、$A,B\neq O$。

\textbf{用词特别注意：}$OA\cdot OB$ 是\emph{线段长度的乘积}，不是向量内积；
$P=(A+B)/2$ 指\emph{坐标中点}。本定理的逻辑方向是
\emph{相切 $\Rightarrow$ 面积为 $ab$}；其逆命题不一定成立。
\end{statementbox}
